%% secoes/04-discussao.tex — incluído por main.tex

%%=====================================================================%%
\section{Discussão}\label{sec:discussao}
%%=====================================================================%%

Esta seção compara os quatro desfechos entre si, relaciona os resultados aos achados da literatura e apresenta as limitações do estudo.

%%---------------------------------------------------------------------%%
\subsection{Padrões comuns e específicos entre os desfechos}\label{subsec:disc-padroes}
%%---------------------------------------------------------------------%%

O número de consultas de pré-natal é o item que mais se repete entre os desfechos, mas com peso muito diferente em cada um. No baixo peso ao nascer, aparece em 23 dos 40 subgrupos e em todos os recortes; no nascimento pré-termo, está entre os cinco primeiros subgrupos dos oito recortes; no Apgar baixo, aparece em três recortes; e na malformação congênita não figura em nenhum subgrupo. Nos três desfechos em que aparece, a faixa de 1 a 3 consultas está associada aos maiores riscos: até 87\% acima da basal no baixo peso, até 103\% no pré-termo e até 86\% no Apgar baixo, quando isolada. A faixa de 4 a 6 consultas também eleva o risco, de 20\% a 46\% acima da basal no baixo peso e de 21\% a 46\% no pré-termo. Como as duas faixas correspondem às categorias do SINASC abaixo de sete ou mais consultas, o resultado indica que o risco permanece elevado mesmo quando o pré-natal foi iniciado, mas não completado.

A primeira gestação é o segundo item mais recorrente. No baixo peso, aparece em seis dos oito recortes e, combinada à faixa de 4 a 6 consultas, forma o subgrupo de maior risco nos municípios de até 10 mil habitantes, na população rural e na quilombola. No pré-termo, a mesma combinação eleva o risco em relação à faixa de 4 a 6 consultas isolada nos três recortes em que aparece, e no Apgar baixo a primeira gestação aparece combinada ao parto normal em três recortes, com risco entre 36\% e 40\% acima da basal.

Outros itens separam os desfechos. O sexo feminino acompanha o baixo peso, enquanto o sexo masculino aparece na malformação congênita, em todos os recortes, e no Apgar baixo, em seis deles. O parto cesáreo está presente nos oito recortes da malformação e no subgrupo de maior risco de seis deles, enquanto o parto normal aparece no Apgar baixo, em quatro recortes, e no baixo peso entre mães menores de 15 anos. Os itens de infraestrutura domiciliar também têm sentidos opostos: no baixo peso e na malformação, descrevem em sua maioria condições presentes, como banheiro, água canalizada e coleta de lixo; no Apgar baixo, descrevem condições ausentes ou precárias, como lixo não coletado, ausência de calçamento, domicílio improvisado e despesa total da família muito baixa.

Entre os segmentos, as mães menores de 15 anos concentram os valores mais altos: têm a maior prevalência basal de baixo peso (11,60\%), de pré-termo (17,26\%) e de malformação (0,63\%), e a segunda maior de Apgar baixo (3,13\%), além do subgrupo de maior risco do baixo peso (2,27 vezes a basal) e do subgrupo de maior risco do Apgar fora da população indígena (2,51 vezes). A população indígena apresenta o menor risco associado à faixa de 1 a 3 consultas, tanto no baixo peso quanto no pré-termo, e é o único segmento em que o nascimento domiciliar aparece como item de risco, no pré-termo (47\% acima da basal) e, sobretudo, no Apgar baixo, em que chega a 4,32 vezes a basal quando combinado a domicílio de 1 a 2 cômodos.

%%---------------------------------------------------------------------%%
\subsection{Relação com a literatura}\label{subsec:disc-literatura}
%%---------------------------------------------------------------------%%

Os itens mais recorrentes no baixo peso ao nascer coincidem com os fatores de maior risco identificados por regressão na mesma coorte: poucas consultas de pré-natal, primiparidade e sexo feminino \cite{falcao2020}. A diferença está no que cada abordagem retorna. A regressão estima o efeito de cada fator ajustado pelos demais; a descoberta de subgrupos mostra em quais combinações o desfecho se concentra e quanto cada combinação representa das ocorrências do segmento. O fato de cinco dos sete segmentos vulneráveis terem prevalência de baixo peso inferior à do recorte geral é compatível com o paradoxo epidemiológico do baixo peso ao nascer descrito para o Brasil, em que as maiores proporções se concentram nas áreas de maior desenvolvimento \cite{silva2010}, mas este estudo não permite separar esse efeito de diferenças de registro entre os segmentos.

No pré-termo, as maiores prevalências entre mães indígenas e mães muito jovens acompanham o padrão nacional descrito por estudos de séries temporais \cite{alberton2023,victor2025}, e o peso do número de consultas é coerente com o observado na Coorte de Nascimentos do CIDACS, em que poucas consultas se associaram tanto à prematuridade incidente quanto à recorrente \cite{rocha2022}. Naquele estudo, a idade materna de 14 a 19 anos e a raça/cor indígena se associaram apenas à prematuridade incidente, o que reforça que o peso de um mesmo atributo varia conforme o subconjunto analisado, como observado aqui entre os segmentos.

No Apgar baixo, os resultados se aproximam dos fatores maternos identificados em âmbito nacional, como menos de seis consultas, baixa escolaridade e raça/cor preta \cite{magalhaes2023}. O nascimento domiciliar entre indígenas, item de maior risco do desfecho, é coerente com o perfil descrito em território indígena do Pará, onde há gravidez na adolescência e baixa proporção de gestantes com sete ou mais consultas \cite{lima2025}. Já a maior prevalência basal entre mães menores de 15 anos difere do estudo nacional que deixou de encontrar a gravidez na adolescência como fator de risco em 2018--2019 \cite{magalhaes2023}; a diferença pode estar no recorte etário, mais restrito aqui, e no período analisado.

Na malformação congênita, os resultados se afastam da literatura. Estudos nacionais associam o desfecho à ausência de pré-natal e à raça/cor preta \cite{fernandes2025,trevilato2022}, enquanto aqui o número de consultas não aparece e os subgrupos de maior risco reúnem parto cesáreo, nascimento hospitalar e condições domiciliares presentes. Uma explicação possível é que esses itens indiquem maior chance de a anomalia ser detectada e registrada, e não maior ocorrência: a prevalência de anomalias no Brasil é considerada afetada pela detecção pré-natal e pela subnotificação, e a região de maior desenvolvimento apresenta as maiores taxas \cite{melkan2025}. Também é possível que o parto cesáreo seja consequência do diagnóstico pré-natal da anomalia, e não fator anterior a ela. Os dados disponíveis não permitem distinguir essas hipóteses.

%%---------------------------------------------------------------------%%
\subsection{Limitações}\label{subsec:disc-limitacoes}
%%---------------------------------------------------------------------%%

Os subgrupos descrevem associações, e não relações causais. Como cada regra é uma conjunção de itens, o risco de um subgrupo não isola o efeito de nenhum item individual, e a comparação entre regras só é direta quando elas diferem por um único item. Alguns atributos, como o tipo de parto e o local do nascimento, são registrados no momento do nascimento e podem ser consequência, e não antecedente, do desfecho ou da condição que o provocou. A WRAcc pondera a cobertura do subgrupo, o que favorece regras amplas, como as definidas apenas pelo sexo do recém-nascido, que aparecem entre as primeiras mesmo com risco baixo. Foram apresentados apenas os subgrupos da primeira execução do SSDP+, e alguns itens relevantes emergiram somente na segunda: no Apgar baixo da população indígena, por exemplo, o nascimento em local classificado como ``outros'' teve \emph{lift} de 6,43, ainda que com apenas 44 casos. Não foram calculados intervalos de confiança para o \emph{lift}.

Nos segmentos menores, as regras abrangem poucos casos, sobretudo na malformação congênita: os subgrupos de maior risco das populações quilombola e de mães menores de 15 anos reúnem 27 e 29 ocorrências, o que torna esses valores mais instáveis; o mesmo vale para os recortes municipais do painel. Algumas regras contêm itens redundantes em relação ao próprio segmento, como a idade materna inferior a 20 anos entre mães menores de 15 anos. Por fim, os dados referem-se a nascimentos de 2011 a 2015 em famílias de baixa renda registradas no Cadastro Único, o que limita a generalização para a população brasileira como um todo; dependem da qualidade de preenchimento do SINASC, em especial para as anomalias congênitas; e os atributos com muitos valores ausentes foram excluídos da formação das regras.
